% Conversão fiel de chapter-05-text-h-0077; origem RELATORIO_ATUAL.odt.
\hypertarget{__RefHeading___Toc4803_3118176633}{}
\refstepcounter{subsection}
\subsection*{5.2.1 Repositório do código Fonte do Projeto}
\label{ages-iii:text-h-0077}
\addcontentsline{toc}{subsection}{5.2.1 Repositório do código Fonte do Projeto}

% SOURCE text-p-1029
\indent Repositório referente à interface de Administrador + Dashboard e Interface de Clientes.\par

% SOURCE text-p-1030
Mais informações podem ser encontradas no repositório \href{https://tools.ages.pucrs.br/sistema-de-controle-inteligente-de-estacionamento-com-rfid/frontend}{RFID Frontend}.\par

% SOURCE text-p-1032
\indent Repositório responsável pela lógica de negócio da aplicação e pela disponibilização dos endpoints consumidos pelas interfaces web. Centraliza autenticação, regras operacionais, persistência de dados e integração com os demais serviços da solução. O projeto foi estruturado para execução em ambiente serverless, com implantação por meio do AWS Lambda\space\cite{ages-iii-aws-lambda}.\par

% SOURCE text-p-1033
Mais informações podem ser encontradas no repositório \href{https://tools.ages.pucrs.br/sistema-de-controle-inteligente-de-estacionamento-com-rfid/backend}{RFID Backend}.\par

% SOURCE text-p-1035
\indent Repositório responsável pela comunicação entre o leitor RFID físico e a aplicação em nuvem. Esse serviço opera no ambiente local onde o equipamento está instalado, utilizando uma API RESTful\space\cite{ages-iii-fielding-rest} para comunicação com o backend, processando os eventos capturados e encaminhando as informações ao backend principal hospedado na nuvem.\par

% SOURCE text-p-1036
Mais informações podem ser encontradas no repositório \href{https://tools.ages.pucrs.br/sistema-de-controle-inteligente-de-estacionamento-com-rfid/gateway}{RFID Gateway}.\par
